\section{VQA on H$_2$ Using STO-3G}

\subsection{Methodology}
We target the electronic ground state of molecular hydrogen at an internuclear separation of 0.735~\r{A} within the STO-3G basis. After mapping the fermionic operators to qubits via the Jordan--Wigner transform, the electronic Hamiltonian acts on a four-qubit Hilbert space and is expressed as a linear combination of Pauli strings:
\begin{equation}
    H_{\text{H}_2} = \sum_{i} c_i\, P_i, \qquad P_i \in \{\mathbb{I}, X, Y, Z\}^{\otimes 4},
\end{equation}
where the real coefficients $c_i$ correspond to the one- and two-electron molecular integrals. The Variational Quantum Eigensolver (VQE) approximates the ground state energy $E_0$ by minimizing the expectation value with respect to a parameterized ansatz $|\psi(\boldsymbol{\theta})\rangle$:
\begin{equation}
    E(\boldsymbol{\theta}) = \langle \psi(\boldsymbol{\theta}) | H_{\text{H}_2} | \psi(\boldsymbol{\theta}) \rangle \ge E_0.
\end{equation}
Our ansatz $|\psi(\boldsymbol{\theta})\rangle = U(\boldsymbol{\theta})|0\rangle^{\otimes 4}$ employs a hardware-efficient architecture consisting of alternating layers of single-qubit $R_y$ rotations and linear CNOT entanglement. We evaluate the performance across circuit depths ranging from three to seven layers, corresponding to a parameter space dimensionality of $N_{\theta} \in \{12, 16, 20, 24, 28\}$. Optimization is performed using the Adam algorithm with a learning rate of $10^{-2}$ and a convergence tolerance of $10^{-7}$, driven by analytic gradients computed via the parameter-shift rule.

\subsection{Performance Analysis}
We benchmark the LogosQ framework against established libraries---Qiskit, PennyLane, and Yao.jl---under identical ansatz configurations and optimizer hyperparameters. LogosQ demonstrates superior computational efficiency, leveraging its compiled Rust backend and native tensor contraction primitives to achieve wall-clock runtimes $2\times$ to $5\times$ faster than the Python-based alternatives (Qiskit and PennyLane) and competitive with the Julia-based Yao.jl.

In terms of solution quality, LogosQ consistently converges to chemical accuracy, maintaining energy errors below $2 \times 10^{-4}$ Hartree across the parameter sweep. Notably, at deeper circuit depths (24--28 parameters), LogosQ achieves near-exact agreement with the full configuration interaction (FCI) energy ($|E - E_{\text{FCI}}| \approx 1 \times 10^{-5}$ Hartree) while requiring fewer optimization iterations to reach convergence compared to other frameworks. This combination of high-performance execution and numerical stability positions LogosQ as a robust platform for intermediate-scale variational quantum algorithms.
